在本书中，我们关注的是神经网络的细节：它们如何工作，以及如何用它们来解决模式识别问题. 这些内容有许多直接的实际应用. 但是，当然，人们对神经网络感兴趣的一个原因是希望有一天它们能远远超越这些基本的模式识别问题. 也许它们，或者某种其他基于数字计算机的方法，最终会被用来建造会思考的机器，即能够匹敌或超越人类智能的机器？这一设想远远超出了本书所讨论的内容——或者说超出了世界上任何人知道如何做到的范围. 但进行推测是很有趣的.

关于计算机是否\emph{可能}匹敌人类智能，一直存在许多争论. 我不打算参与这个问题. 尽管争论不断，我相信智能计算机是可能的这一点并不存在严重疑问——尽管它可能极其复杂，也许远远超出当前技术——而当前的否定者有一天会看起来很像活力论者.

更确切地说，我在这里探讨的问题是：是否存在一套\emph{简单}的原则可以用来解释智能？具体而言，更实际地说，是否存在一个\emph{智能的简单算法}？

存在真正简单的智能算法这一想法是一个大胆的想法. 它听起来可能过于乐观而不像是真的. 许多人有一种强烈的直觉，认为智能具有相当大的不可约复杂性. 他们对人类思维的惊人多样性和灵活性印象深刻，以至于得出结论：智能的简单算法必定是不可能的. 尽管有这种直觉，我认为匆忙下结论是不明智的. 科学史上充满了这样的例子：一个现象最初看起来极其复杂，但后来被一些简单而有力的思想所解释.

例如，考虑天文学的早期. 人类自古以来就知道天空中有各种各样的物体：太阳、月亮、行星、彗星和恒星. 这些物体的运动方式非常不同——例如，恒星以庄重、规律的方式在天空中移动，而彗星则仿佛凭空出现，划过天空，然后消失. 在 16 世纪，只有愚蠢的乐观主义者才会想象所有这些物体的运动可以用一套简单的原则来解释. 但在 17 世纪，Newton 提出了他的万有引力理论，它不仅解释了所有这些运动，还解释了潮汐和地球上抛射体行为等地面现象.16 世纪的愚蠢乐观主义者回过头来看倒像是悲观主义者，要求的太少了.

当然，科学中还有更多这样的例子. 考虑构成我们世界的无数化学物质，Mendeleev 的元素周期表如此优美地解释了它们，而周期表本身又由几条可以从量子力学中得到的简单规则所解释. 或者生物世界中为何有如此多的复杂性和多样性这一谜题，其起源原来在于自然选择的进化原理. 这些以及许多其他例子表明，仅仅因为我们的头脑——目前智能的最佳范例——所做的\emph{看起来}非常复杂，就排除对智能的简单解释，是不明智的*\footnote{在本附录中，我假设一台计算机要被认定为智能的，其能力必须匹敌或超越人类思维能力. 因此我将问题``存在智能的简单算法吗？''等同于``是否存在一个能够沿着与人类大脑基本相同的思路进行`思考'的简单算法？''然而，值得注意的是，很可能存在一些不包含人类思维、但以有趣的方式超越人类思维的智能形式.}.

相反地，尽管有这些乐观的例子，逻辑上也有可能智能只能由大量根本不同的机制来解释. 就我们的大脑而言，那些许多机制也许是在我们物种的进化历史中响应许多不同的选择压力而进化出来的. 如果这种观点是正确的，那么智能涉及相当大的不可约复杂性，智能的简单算法就是不可能的.

这两种观点哪一种正确？

为了深入了解这个问题，让我们问一个密切相关的问题，即是否存在关于人类大脑如何工作的简单解释. 具体来说，让我们看看一些量化大脑复杂性的方法. 我们的第一种方法是从连接组学看大脑. 这完全关乎原始的布线：大脑中有多少神经元，多少胶质细胞，以及神经元之间有多少连接. 你以前可能听说过这些数字——大脑中大约有 1000 亿个神经元，1000 亿个胶质细胞，以及神经元之间 100 万亿个连接. 这些数字令人震惊. 它们也令人生畏. 如果我们需要理解所有这些连接的细节（更不用说神经元和胶质细胞了）才能理解大脑如何工作，那么我们肯定不会得到一个智能的简单算法.

还有第二种更乐观的观点，即从分子生物学看大脑. 其思路是问需要多少遗传信息来描述大脑的结构. 为了把握这个问题，我们将从考虑人类和黑猩猩之间的遗传差异开始. 你可能听说过``人类有 98\% 是黑猩猩''这样的说法. 这种说法有时会有变化——流行的变体也给出 95\% 或 99\% 这样的数字. 这些变化之所以出现，是因为这些数字最初是通过比较人类和黑猩猩基因组的样本而非整个基因组来估计的. 然而，在 2007 年，整个黑猩猩基因组被测序了（另见此处），我们现在知道人类和黑猩猩的 DNA 在大约 1.25 亿个 DNA 碱基对上存在差异. 而每个基因组中总共有大约 30 亿个 DNA 碱基对. 所以说人类有 98\% 是黑猩猩是不对的——我们更像是 96\% 是黑猩猩.

那 1.25 亿个碱基对中有多少信息？每个碱基对可以用四种可能性之一来标记——遗传密码的``字母''，即碱基腺嘌呤、胞嘧啶、鸟嘌呤和胸腺嘧啶. 因此每个碱基对可以用两位信息来描述——刚好足够指定四个标签之一. 所以 1.25 亿个碱基对相当于 2.5 亿位信息. 这就是人类和黑猩猩之间的遗传差异！

当然，那 2.5 亿位信息包含了人类和黑猩猩之间的所有遗传差异. 我们只对与大脑相关的差异感兴趣. 不幸的是，没有人知道总遗传差异中有多大比例是用来解释大脑差异的. 但为了论证起见，让我们假设那 2.5 亿位信息中大约有一半用于大脑差异. 那就是总共 1.25 亿位信息.

1.25 亿位是一个令人印象深刻的大数字. 让我们通过将其转化为更人性化的术语来感受它有多大. 具体来说，相当于多少英文文本？事实证明，英文文本的信息含量大约是每字母 1 位. 这听起来很低——毕竟字母表有 26 个字母——但英文文本中存在大量的冗余. 当然，你可能会争辩说我们的基因组也是冗余的，所以每个碱基对两位是高估了. 但我们将忽略这一点，因为最坏的情况下这意味着我们高估了大脑的遗传复杂性. 在这些假设下，我们看到我们的大脑和黑猩猩大脑之间的遗传差异大约相当于 1.25 亿个字母，或大约 2500 万个英文单词. 这大约是钦定版圣经的 30 倍.

那是很多信息. 但并非大到不可理解. 它处于人类的尺度上. 也许没有哪个人能理解那代码中所写的全部内容，但一群人也许可以通过适当的专业化来集体理解它. 而且尽管信息量很大，与描述我们大脑中 1000 亿个神经元、1000 亿个胶质细胞和 100 万亿个连接所需的信息相比，它就微不足道了. 即使我们使用一个简单、粗略的描述——比如说，用 10 个浮点数来表征每个连接——那也需要大约 70 千万亿位. 这意味着遗传描述比人类大脑的完整连接组在复杂性上大约少了五亿倍.

我们从中了解到的是，我们的基因组不可能包含所有神经连接的详细描述. 相反，它必定只指定了大脑的总体结构和基本原理. 但那个结构和那些原理似乎足以保证我们人类会成长为智能的. 当然，有一些注意事项——成长中的儿童需要健康、有刺激的环境和良好的营养来实现其智力潜能. 但只要我们在一个合理的环境中成长，一个健康的人将拥有非凡的智能. 在某种意义上，我们基因中的信息包含了我们如何思考的精髓. 而且此外，那遗传信息中包含的原理似乎很可能在我们集体把握的能力范围之内.

上面所有的数字都是非常粗略的估计. 有可能 1.25 亿位是一个极大的高估，存在一套更紧凑得多的核心原理支撑着人类思维. 也许那 1.25 亿位中的大部分只是对相对次要细节的微调. 或者也许我们在计算这些数字时过于保守了. 显然，如果那是真的，那就太好了！对于我们当前的目的，关键点是：大脑的结构是复杂的，但它远没有你根据大脑中连接数量所想象的那么复杂. 从分子生物学看大脑的观点表明，我们人类有一天应该能够理解大脑结构背后的基本原理.

在最后几段中，我忽略了 1.25 亿位仅仅量化了人类和黑猩猩大脑之间的遗传\emph{差异}这一事实. 并非我们所有的脑功能都归因于那 1.25 亿位. 黑猩猩本身就是非凡的思考者. 也许智能的关键主要在于黑猩猩和人类共有的心智能力（和遗传信息）. 如果这是正确的，那么人类大脑可能只是黑猩猩大脑的一个小升级，至少就底层原理的复杂性而言. 尽管传统的人类沙文主义关于我们独特能力的说法，这并非不可想象：黑猩猩和人类的遗传谱系仅仅在 500 万年前分化，这在进化时间尺度上只是一眨眼. 然而，在缺乏更有说服力的论据的情况下，我同情传统的人类沙文主义：我的猜测是，支撑人类思维的最有趣的原理在于那 1.25 亿位中，而不在于我们与黑猩猩共享的那部分基因组中.

采用从分子生物学看大脑的观点使我们描述的复杂性减少了大约九个数量级. 虽然令人鼓舞，但它并没有告诉我们智能的简单算法是否可能. 我们能否进一步降低复杂性？而且，更关键的是，我们能否解决智能的简单算法是否可能这个问题？

不幸的是，目前还没有足够强的证据来决定性地解决这个问题. 让我描述一些可用的证据，但需要说明的是，这是一个非常简短且不完整的概述，旨在传达一些近期工作的风格，而非全面综述已知的内容.

在暗示可能存在智能的简单算法的证据中，有一项 2000 年 4 月发表在\emph{Nature}杂志上的实验. 由 Mriganka Sur 领导的一个科学家团队``重新布线''了新生雪貂的大脑. 通常，雪貂眼睛的信号被传输到大脑中称为视觉皮层的部分. 但对于这些雪貂，科学家们取出来自眼睛的信号并重新路由，使其改为传到听觉皮层，即通常用于听觉的大脑区域.

为了理解他们这样做时发生了什么，我们需要了解一点关于视觉皮层的知识. 视觉皮层包含许多方向柱. 这些是神经元的小板块，每个都对来自某个特定方向的视觉刺激做出反应. 你可以把方向柱想象成微小的方向传感器：当有人从某个特定方向照射亮光时，一个相应的方向柱被激活. 如果光被移动，一个不同的方向柱被激活. 视觉皮层中最重要的高级结构之一是方向图，它描绘了方向柱是如何布局的.

科学家们发现，当来自雪貂眼睛的视觉信号被重新路由到听觉皮层时，听觉皮层发生了变化. 方向柱和方向图开始在听觉皮层中出现. 它比通常在视觉皮层中发现的方向图更无序，但无疑相似. 此外，科学家们做了一些简单的测试，测试雪貂如何对视觉刺激做出反应，训练它们在光从不同方向闪烁时做出不同的反应. 这些测试表明，雪貂仍然可以学会``看见''，至少以基本的方式，使用听觉皮层.

这是一个惊人的结果. 它表明存在共同的原则支撑着大脑不同部分如何学会对感觉数据做出反应. 这种共同性至少为存在一套支撑智能的简单原则这一想法提供了一些支持. 然而，我们不应该自欺欺人地认为这些实验中雪貂的视力有多好. 行为测试只测试了视觉的非常粗略的方面. 而且，当然，我们无法问雪貂它们是否``学会了看见''. 所以这些实验并没有证明重新布线的听觉皮层给了雪貂高保真的视觉体验. 因此它们只提供了有限的证据来支持共同原则支撑大脑不同部分如何学习这一想法.

有什么证据反对智能的简单算法这一想法？一些证据来自进化心理学和神经解剖学领域. 自 1960 年代以来，进化心理学家发现了广泛的\emph{人类共性}，即所有人类共有的复杂行为，跨越文化和教养. 这些人类共性包括母子之间的乱伦禁忌、音乐和舞蹈的使用，以及许多复杂的语言结构，如脏话（即禁忌词）、代词的使用，甚至像动词这样基本的结构. 作为这些结果的补充，大量来自神经解剖学的证据表明，许多人类行为由大脑特定的局部区域控制，而这些区域在所有人中似乎都是相似的. 综合来看，这些发现表明许多非常专门化的行为被硬连线到我们大脑的特定部分.
有些人从这些结果中得出结论，认为这许多大脑功能必定需要各自不同的解释，因此大脑的功能具有一种不可约化的复杂性，这种复杂性使得对大脑运作的简单解释（也许还有智能的简单算法）成为不可能. 例如，持这种观点的一位著名人工智能研究者是 Marvin Minsky. 在 20 世纪 70 年代和 80 年代，Minsky 发展了他的``心智社会''理论，其基础观念是：人类智能是大量各自简单（但非常不同）的计算过程所组成的一个庞大社会的结果，Minsky 将这些过程称为智能体. 在描述这一理论的书中，Minsky 总结了他所认为的这一观点的力量：

\begin{quote}
是什么神奇的诀窍让我们变得智能？诀窍就是没有诀窍. 智能的力量源于我们巨大的多样性，而非任何单一的、完美的原理.
\end{quote}

在回应* *见 ``Contemplating Minds: A Forum for Artificial Intelligence''，由 William J. Clancey、Stephen W. Smoliar 和 Mark Stefik 编辑（MIT Press，1994）. 对其著作的评论时，Minsky 详细阐述了心智社会的动机，给出了一个与上述论证类似的论证，其基础是神经解剖学和进化心理学：

\begin{quote}
我们现在知道，大脑本身由数百个不同的区域和核团组成，每一个都具有显著不同的结构要素和排列方式，而且其中许多区域和核团参与了我们心理活动中明显不同的方面. 这一现代知识的大量积累表明，许多传统上用``智能''或``理解''这类常识术语来描述的现象，实际上涉及复杂的机器组合.
\end{quote}

当然，Minsky 并不是唯一持有这类观点的人；我只是把他作为这一论证路线支持者的一个例子. 我觉得这个论证很有趣，但不认为其证据令人信服. 虽然大脑确实由大量具有不同功能的不同区域组成，但并不能因此就推断出对大脑功能的简单解释是不可能的. 也许那些结构上的差异源于共同的基本原理，正如彗星、行星、太阳和恒星的运动都源于单一的引力一样. 无论是 Minsky 还是其他任何人，都没有令人信服地论证反对这样的基本原理.

我自己的偏见是倾向于存在一种智能的简单算法. 我喜欢这个想法的主要理由，除了上述（不具决定性的）论证之外，还在于它是一个乐观的想法. 在研究工作中，一种没有根据的乐观往往比一种看似更有根据的悲观更有成效，因为乐观者有勇气出发去尝试新事物. 那是通往发现的道路，即使所发现的东西也许并非最初所希望的. 悲观者在某种狭隘的意义上也许更``正确''，但会比乐观者发现更少的东西.

这种观点与我们通常评判想法的方式形成了鲜明对比：我们通常试图弄清楚它们是对还是错. 对于处理日常研究中的琐碎细节来说，这是一种明智的策略. 但对于评判一个宏大而大胆的想法——那种定义一个完整研究纲领的想法——这可能是错误的方式. 有时候，我们只有微弱的证据来判断这样一个想法是否正确. 我们可以温顺地拒绝追随这个想法，转而把所有时间都花在眯着眼审视现有证据上，试图辨别什么是真的. 或者我们可以接受目前还没有人知道，转而努力发展这个宏大而大胆的想法，因为我们明白，虽然我们没有成功的保证，但只有这样才能推进我们的理解.

话虽如此，在其\emph{最}乐观的形式下，我不相信我们会找到一种智能的简单算法. 更具体地说，我不相信我们会找到一个真正简短的 Python（或 C 或 Lisp，或别的什么）程序——比方说，最多一千行代码——来实现人工智能. 我也不认为我们会找到一个真正容易描述的、能够实现人工智能的神经网络. 但我确实相信，值得像我们能够找到这样一个程序或网络那样去行动. 那是通往洞见的道路，沿着这条道路前进，我们也许有一天会理解得足够多，从而写出一个更长的程序或构建一个更复杂的网络，它确实展现出智能. 因此，值得像存在一种极其简单的智能算法那样去行动.

在 20 世纪 80 年代，著名数学家和计算机科学家 Jack Schwartz 受邀参加一场人工智能支持者与人工智能怀疑者之间的辩论. 辩论变得失控，支持者对即将到来的惊人事物做出了夸张的断言，而怀疑者则加倍坚持他们的悲观态度，声称人工智能根本不可能.Schwartz 是这场辩论的局外人，在讨论升温时保持沉默. 在一个间歇，有人请他发言，陈述他对所讨论问题的看法. 他说：``嗯，其中一些进展可能还差一百个诺贝尔奖那么远''（ref，第 22 页）. 在我看来，这是一个完美的回应. 人工智能的关键在于简单而强大的想法，我们可以而且应该乐观地寻找那些想法. 但我们将需要许多这样的想法，而且我们还有很长的路要走！
